\documentclass[11pt,a4paper]{book}
\usepackage[a4paper,margin=2.6cm]{geometry}
\usepackage{amsmath,amssymb}
\usepackage{graphicx}
\usepackage{booktabs}
\usepackage{longtable}
\usepackage{enumitem}
\usepackage{titlesec}
\usepackage{hyperref}
\usepackage{fancyhdr}
\usepackage[T1]{fontenc}
\usepackage{lmodern}
\usepackage{xcolor}

\hypersetup{colorlinks=true,linkcolor=blue!60!black,citecolor=blue!60!black,urlcolor=blue!60!black}
\titleformat{\chapter}[display]{\normalfont\huge\bfseries}{\chaptertitlename\ \thechapter}{20pt}{\Huge}
\pagestyle{fancy}
\fancyhf{}
\fancyhead[LE,RO]{\thepage}
\fancyhead[RE]{\nouppercase{\leftmark}}
\fancyhead[LO]{\nouppercase{\rightmark}}
\renewcommand{\headrulewidth}{0.4pt}

\title{\textbf{Lunar-V2 Thermal Pipeline}\\[0.4em]\Large A Four-Phase Roadmap from Point Validation to TSUKIMI Retrieval}
\author{R.\,P.\,Gregorio\\[0.3em]\small TSUKIMI Mission Thermal-Model Support}
\date{\today}

\begin{document}

\frontmatter
\maketitle

\chapter*{Scope}
\addcontentsline{toc}{chapter}{Scope}
This document is the living design pipeline for the Lunar-V2 thermal
modelling codebase. It describes, phase by phase, how the repository
evolves from a single-point Apollo HFE validation into a retrieval-ready
forward operator suitable for TSUKIMI mission science. Each chapter
specifies the scope, the papers it rests on, the data sources, the
numerics it adds or replaces, and the quantitative success criterion by
which progress is measured. The pipeline is \emph{cumulative}: every
later phase uses the solver and infrastructure locked in earlier
phases.

\tableofcontents

\mainmatter

\chapter{Pipeline Overview}
\section{Motivation}
The TSUKIMI mission requires a forward model mapping regolith
physical state $(H,\rho,\text{rock fraction},\phi_\text{ice})$ to
emerging thermal-infrared brightness temperatures at the instrument's
observation geometry. The retrieval chain is only as accurate as the
weakest link in that forward operator. Lunar-V2 therefore starts by
benchmarking the unmodified Hayne et al.\ (2017) thermal model at the
only in-situ lunar subsurface truth available --- the Apollo 15 and 17
Heat Flow Experiments --- and then augments that model in three
successive phases until it reaches the fidelity the mission needs.

\section{Phases at a glance}
\begin{longtable}{p{0.16\textwidth} p{0.32\textwidth} p{0.48\textwidth}}
\toprule
\textbf{Phase} & \textbf{Headline deliverable} & \textbf{Why it matters} \\
\midrule
\textbf{1} --- Point validation & Hayne 2017 vs Discrete Layer vs Apollo 15/17 HFE stabilised-window deep sensors ($z\!\ge\!80$~cm). Reproduction of Hayne 2017 Figure~4. & Establishes the absolute accuracy envelope of the forward operator at one point where ground truth exists. \\
\textbf{2} --- Improved global Hayne & One merged global model combining (i) latitude-dependent H-parameter map, (ii) Bandfield 2011 rock-abundance mixing, (iii) Martinez \& Siegler 2021 cold-region conductivity, (iv) B\"urger 2024 microphysical scaling, (v) LOLA-DEM horizon + shadow mask. Validated against Diviner PCP global maps. & Extends the one-point model to the full Moon, in particular to PSRs where the Hayne $K(T,z)$ form breaks down and shadowing dominates. \\
\textbf{3} --- RTM coupling, Jacobian, ice-stability & Two-stream regolith radiative-transfer coupling; finite-difference Jacobian of brightness-$T$ w.r.t.~retrieval state; ice-stability index map (Schorghofer 2008). & Turns the forward model into a TSUKIMI-ready retrieval operator. \\
\textbf{4} --- Thesis integration + future work & LaTeX thesis chapters, published figure bundle, scoping of post-thesis extensions (small-scale roughness, time-dependent ice retreat, Mars adaptation, in-flight retrieval harness). & Completes the dissertation and delimits what is in vs.\ out of scope. \\
\bottomrule
\end{longtable}

\section{Success definition}
Lunar-V2 is considered \emph{pipeline complete} when, at the end of
Phase 3, the Jacobian of the forward model has condition number
$\le 10^{4}$ on the TSUKIMI retrieval state vector over a
polar-stereographic grid, and the forward model's bolometric residual
against Diviner PCP has RMS $\le 6$~K at one pole (beating the
published Hayne 2017 $\sim 8$~K envelope).

%----------------------------------------------------------------------
\chapter{Phase 1 --- Point-Source Apollo Validation}
\section{Scope}
Phase 1 is \emph{strictly} one-point-at-a-time. It uses the Hayne 2017
forward solver as shipped in \texttt{phayne/heat1d} (vendored under
\texttt{third\_party/heat1d/}) and compares it, and the Apollo-
calibrated Discrete Layer alternative, to the Apollo 15 and 17 HFE
stabilised-window deep sensors.

\section{Model}
The solver is the 1-D Crank-Nicolson finite-volume integrator with
geometric depth grid. Properties exactly match Hayne's GitHub:
\begin{align}
K(T,z) &= K_c(z)\bigl[1 + \chi(T/350)^3\bigr],
&& K_c(z) = K_d - (K_d-K_s)\,e^{-z/H}, \\
\rho(z) &= \rho_d - (\rho_d-\rho_s)\,e^{-z/H}, \\
A(i) &= A_0 + a\left(i/45^\circ\right)^3 + b\left(i/90^\circ\right)^8,
&& (A_0,a,b) = (0.12,0.06,0.25),
\end{align}
with $c_p(T)$ the Hayne 2017 Appendix~A fourth-order polynomial.
Boundary conditions are radiative at the surface and geothermal flux
($Q_b = 21$~mW\,m$^{-2}$ for Apollo 15, Langseth 1976) at the base.

\section{Validation targets}
\begin{itemize}
\item Apollo 15 deep-sensor RMSE $\le 1$~K (Discrete Layer) and $\le 1.5$~K (Hayne 2017).
\item Apollo 17 mean-$T$ at $z = 1$~m within $\pm 2$~K of observed.
\item Hayne 2017 Figure~4 reproduced at the Apollo 15 site: $T(z)$ at
      four local times, overlaid with HFE stabilised-window deep sensors.
\item Summary statistics written to \texttt{output/tables/apollo\_stats.csv}.
\end{itemize}

\section{What is NOT in Phase 1}
\begin{itemize}
\item No Chang'E-4 or ChaSTE probe comparisons (raw time-series are login-gated and add no information beyond Apollo for a one-point benchmark).
\item No global maps, no latitude-dependent H-parameter, no rock abundance, no cold-region conductivity correction.
\item The DEM/horizon/shadowing infrastructure exists in the code base but is not active for the Apollo flat-site runs.
\end{itemize}

%----------------------------------------------------------------------
\chapter{Phase 2 --- Improved Hayne 2017 Global Model}
\section{Motivation}
The Hayne 2017 GitHub code implements only the per-point solver. Its
global maps (Figure~7, H-parameter latitude fit, rock-abundance
mixing, bolometric emissivity) were analysis products that never made
it into the committed codebase. Phase 2 re-builds those pieces and,
additionally, incorporates improvements published after the 2017 paper.

\section{Scope}
Phase 2 produces \emph{one} merged forward model with five
enhancements:
\begin{enumerate}[label=(\roman*)]
\item H-parameter latitude map (Hayne 2017 \S 5.2, Eq.~A.9) fitted via \texttt{pyshtools}.
\item Rock-abundance anisothermal two-component temperature field (Bandfield et al.\ 2011).
\item Martinez \& Siegler (2021) density-dependent conductivity in the cold regime where the Hayne $K(T,z)$ form breaks down below $\sim 80$~K.
\item B\"urger et al.\ (2024, \emph{JGR Planets}) microphysical thermal model --- adds the latitudinal dependence of radiative contact-conductivity coupling.
\item DEM + horizon + shadow-mask + sky-view-factor pipeline from LOLA 80~m/pixel polar DEM.
\end{enumerate}

\section{Validation}
Compared against Diviner Polar Cumulative Products (Williams
et al.\ 2019, 2020), which are auto-fetched from the NASA PDS Geosciences
Node via \texttt{boot.ensure\_diviner\_pcp()}. The target is
\textbf{bolometric-$T$ residual RMS $\le 6$~K at one pole} --- i.e.
beating the published Hayne 2017 $\sim 8$~K envelope.

\section{Success criterion}
A side-by-side figure where, at every polar-stereographic pixel, the
Phase-2 model's bolometric temperature matches Diviner to within the
$\sim 6$~K envelope, and a table quantifying where each of the five
enhancements contributes residual reduction.

%----------------------------------------------------------------------
\chapter{Phase 3 --- RTM Coupling, Jacobian, Ice-Stability}
\section{Scope}
With the improved global forward model locked, Phase 3 turns it into a
retrieval operator.

\paragraph{Radiative transfer.} Couple the Phase-2 surface-$T$ and
$T(z)$ fields to a two-stream regolith RT model to produce emerging
brightness-$T$ spectra in the Diviner channels and at the TSUKIMI
thermal band.

\paragraph{Jacobian.} Compute the finite-difference Jacobian of the
emerging brightness-$T$ with respect to the retrieval state vector
$(H, \text{rock frac}, \rho_s, \phi_\text{ice})$. Cross-check against
an adjoint-state estimate for numerical stability.

\paragraph{Ice stability.} Derive an H$_2$O/CO$_2$ ice-stability depth
map (Schorghofer 2008 criterion) with uncertainty bars propagated from
Phase-2 forward-model residuals.

\section{Success criterion}
Jacobian condition number $\le 10^{4}$ on the TSUKIMI retrieval state
vector across the polar-stereographic grid. Ice-stability map agrees
with Schorghofer 2008's published extents at the 90\% level.

%----------------------------------------------------------------------
\chapter{Phase 4 --- Thesis Integration and Future Work}
\section{Thesis integration}
All figures and tables produced in Phases~1--3 are regenerated by a
single pair of notebooks (\texttt{12\_thesis\_figure\_bundle.ipynb}
and \texttt{13\_paper\_tables.ipynb}) that read cached solver outputs
and emit 300~DPI PDFs + LaTeX tables directly into the thesis source
tree. \texttt{make thesis} produces the manuscript PDF.

\section{Candidate future work (non-thesis)}
\begin{enumerate}
\item \textbf{Small-scale surface roughness} as an additional forward-model layer (Bandfield 2015, Warren 2019) --- shrinks residual brightness-$T$ bias at high-phase-angle retrievals.
\item \textbf{Time-dependent ice retreat.} Couple the Phase-3 ice-stability map to a Diviner-constrained surface-age model and forward-model ice-loss rates.
\item \textbf{Mars adaptation.} Port the improved solver to MGS-MOLA + Mars-Climate-Database for a cross-target validation.
\item \textbf{TSUKIMI in-flight retrieval harness.} Operational wrapper turning the Jacobian into a real Levenberg-Marquardt retrieval with the radiometric-calibration observation-error covariance.
\end{enumerate}

\section{Deliverables}
\begin{itemize}
\item Finalised thesis PDF.
\item Companion manuscript submission (PSJ / JGR / Icarus).
\item Mission-deliverable retrieval operator + Jacobian.
\item This pipeline document, updated to reflect what was actually built.
\end{itemize}

%----------------------------------------------------------------------
\appendix
\chapter{Repository layout}
\begin{verbatim}
Lunar-V2/
├── lunar/                   Python package (solver, grid, properties, I/O)
│   ├── _bootstrap.py        Auto-download and environment setup
│   ├── constants.py         Every numerical constant, cited
│   ├── solver.py            1-D Crank-Nicolson finite-volume integrator
│   ├── properties.py        K(T,z), rho(z), c_p(T), albedo_angle(i)
│   ├── validation.py        Apollo HFE + Diviner PCP loaders
│   └── plotting/            Publication-quality style guide
├── notebooks/
│   ├── phase0_quickstart/   Library smoke test (00_quickstart)
│   ├── phase1_validation/   Apollo 15/17 + Hayne Fig 4 reproduction
│   ├── phase2_improved_hayne/  Global model upgrade (in progress)
│   ├── phase3_rtm_ice/      RTM + Jacobian + ice stability (planned)
│   └── phase4_thesis/       Thesis figure/table bundling (planned)
├── third_party/
│   └── heat1d/              Vendored Paul Hayne heat1d GitHub repo
├── paper/
│   ├── pipeline/            This document
│   └── thesis/              Thesis source (Phase 4)
├── archive/                 Retired scripts + legacy docs (committed)
├── data/                    External datasets (git-ignored, auto-download)
└── output/                  Figures, gifs, tables (git-ignored)
\end{verbatim}

\chapter{References of note}
\begin{itemize}
\item Hayne et al.\ (2017) Global Regolith Thermophysical Properties of the Moon from the Diviner Lunar Radiometer Experiment. \emph{JGR Planets}, 122, 2371--2400.
\item Vasavada, Bandfield, Greenhagen et al.\ (2012) Lunar equatorial thermal environment from Diviner. \emph{JGR Planets}, 117, E00H18.
\item Bandfield et al.\ (2011) Lunar surface rock abundance and regolith fines temperatures from Diviner. \emph{JGR Planets}, 116, E00H02.
\item Martinez \& Siegler (2021) A global thermal conductivity model for lunar regolith at low temperatures. \emph{JGR Planets}, 126.
\item B\"urger et al.\ (2024) A Microphysical Thermal Model for the Lunar Regolith. \emph{JGR Planets}.
\item Schorghofer (2008) The lifetime of ice on main belt asteroids. \emph{Ap.\,J.}, 682, 697--705.
\item Nagihara et al.\ (2018) Examination of the long-term subsurface warming observed at the Apollo 15 and 17 sites. \emph{JGR Planets}, 123.
\item Williams et al.\ (2019, 2020) The global surface temperatures of the Moon as measured by the Diviner Lunar Radiometer Experiment. \emph{Icarus} (2019), \emph{JGR Planets} (2020).
\end{itemize}

\end{document}
